\clearpage
\fancyhead[L]{\small PIAZZI. STERNVERZEICHNISS.}
\hypertarget{native-595-lower}{}
\begin{center}\small Göttingische gelehrte Anzeigen. Stück 91. Seite 903..904. 1815 Juni 10.\\[3pt]\rule{.72\linewidth}{.3pt}\end{center}

\emph{Praecipuarum stellarum inerrantium positiones mediae ineunte saeculo XIX ex observationibus habitis in specula Panormitana ab anno} 1792 \emph{ad annum} 1813. \emph{Ex regia typographia militari.} Palermo 1814. 187 Seiten in Folio. (Preis, in Florenz bei Molini, Landi und Co., 30 Lire.)

Seit der Erscheinung des berühmten \textsc{Piazzi}'schen Sternverzeichnisses von 1803, der Frucht einer ununterbrochenen zehnjährigen angestrengten Arbeit, waren die Fixsternbeobachtungen auf der Palermer Sternwarte theils von \textsc{Piazzi} selbst, theils von seinem geschickten Gehülfen \textsc{Nicolao Cacciatore}, beständig fortgesetzt, erweitert und selbst neu begründet, indem die Fundamentalrectascensionen von $\alpha$ im Adler und $\alpha$ im kleinen Hund durch unmittelbare Vergleichung mit der Sonne fest gestellt wurden. Der vorliegende neue Catalog liefert nun die vollständige Ausbeute der ganzen zwanzigjährigen Arbeit, die Stellungen von 7646 Fixsternen mit aller der Schärfe bestimmt, welche der heutigen beobachtenden Astronomie erreichbar ist.

Eine besondere Sorgfalt hat \textsc{Piazzi} in diesem classischen Werke auch der eignen Bewegung der Sterne gewidmet. Er hat dabei vorzüglich die Beobachtungen von \textsc{Mayer} und \textsc{Bradley} zum Grunde gelegt, sehr häufig auch die Vergleichung seiner eignen Beobachtungen unter sich mit zu Rathe gezo\SourceLeaf{596}gen. Die Stellungen von 1041 Sternen für 1756, wie sie ein anderer Gehülfe \textsc{Piazzi}'s \textsc{Joseph Pilati} zu diesem Behuf aus \textsc{Bradley}'s Beobachtungen reducirt hat, sind am Schlusse des Werks beigefügt.

Die jährliche Präcession der Sterne wurde nach folgenden Formeln berechnet:
\begin{center}\small\renewcommand{\arraystretch}{1.18}\begin{tabular}{@{}l@{\hspace{8pt}}l@{}}
in gerader Aufsteigung $.\;.\;.\;.\;.\;.$ & $46''0395+20''0642\sin\alpha\operatorname{tang}\delta$ \\
in Declination $.\;.\;.\;.\;.\;.\;.\;.\;.\;.$ & $20''0642\cos\alpha$ \\
\end{tabular}\end{center}
welche Zahlen die jährliche Lunisolarpräcession in der Länge $50''388$, die jährliche Bewegung der Aequinoctialpunkte auf dem Aequator vermöge der Planetarischen Einwirkungen $0''1814$ und die mittlere Schiefe der Ekliptik für 1800, $23^\circ\;27'\;55''5$ zum Grunde liegen. Diese Bestimmungen kommen sehr nahe mit denen von \textsc{Bessel} überein. (M. s. das 19. Stück unsrer Blätter 1814 [S. 561 d. B.]).

Das Aeussere des Werks ist zwar des innern Werthes nicht unwürdig, jedoch der Druck etwas öconomischer eingerichtet als bei dem Verzeichniss von 1803. Die scharfe Angabe der geraden Aufsteigungen ist blos in Bogen angesetzt, die in Zeit blos auf Minuten. Alle Sterne sind nach ihren geraden Aufsteigungen geordnet, in 24 Stunden abgetheilt, und in jeder einzelnen Stunde mit fortlaufenden Zahlen bezeichnet. Bei allen Sternen, die besondere Namen haben, sind diese vorzüglich aus \textsc{Ulugh Beighs} Verzeichniss beigefügt, dagegen aber bei solchen Sternen alle weitere Bezeichnung weggeblieben. Dies letztere will uns nicht gefallen, und wir glauben, dass alle Astronomen es unbequem finden werden, statt der ihnen geläufigen \textsc{Bayer}'schen Buchstaben und \textsc{Flamstead}'schen Zahlen nur die ihnen fast sämmtlich fremden Arabischen Namen anzutreffen.

Am Schluss jeder einzelnen Stunde finden sich immer reichhaltige Anmerkungen über eine Menge der darin vorkommenden Sterne.

Nachdem der Catalog schon ganz vollendet und selbst abgedruckt war, machte \textsc{Piazzi} noch die Vergleichung mit den Verzeichnissen der Zodiakalsterne von \textsc{Zach} und \textsc{Barry}'s; die Uebereinstimmung mit dem erstern war fast durchgehends sehr gross; hingegen fanden sich bei den Mannheimer Declinationen häufigere und bedeutendere Unterschiede, wovon indessen nur die grössten in ganze Minuten gehenden hier angeführt sind.

\clearpage
\fancyhead[L]{\small MÜLLER. BERECHNUNG DER SONNENFINSTERNISSE.}
\hypertarget{native-596-lower}{}
\begin{center}\small Göttingische gelehrte Anzeigen. Stück 20. Seite 200. 1816 Februar 3.\\[3pt]\rule{.72\linewidth}{.3pt}\end{center}

\emph{Ausführliche Anleitung zur trigonometrischen Berechnung der an einem gegebenen Ort der Erdfläche sichtbaren Sonnenfinsternisse}, nach zwei verschiedenen sehr genauen Methoden, erläutert durch die Bestimmungen der Erscheinungen der grossen Sonnenfinsterniss des 19. Novembers 1816 für den Nürnberger Meridian, von \textsc{Johann Wolfgang Müller}, Professor der Mathematik am Königl. Gymnasium zu Nürnberg. Sulzbach. In des Commerzienraths J. E. \textsc{Seidel} Kunst- und Buchhandlung. 128 Seiten in Octav, nebst einer Kupfertafel.

Der Verfasser hat die Hauptphasen der merkwürdiegen Sonnenfinsterniss des gegenwärtigen Jahrs für den Nürnberger Horizont, wo ihre Grösse 10 Zoll 38 Minuten betragen wird, mit vieler Sorgfalt nach der trigonometrischen Methode des Neunzigsten im Voraus berechnet, einmal ganz nach \textsc{Tob. Mayer}'s Anordnung, und zweitens nach einem etwas, obwohl eigentlich nicht wesentlich, abgeänderten Gange der Rechnung, und die sämmtlichen Operationen bis ins kleinste Detail in vorliegender kleinen Schrift abdrucken lassen. Anfängern und ungeübten Liebhabern, welche gern eine ähnliche Rechnung für ihren Wohnort ausführen möchten, kann sie als eine brauchbare und fassliche Vorschrift empfohlen werden.
